\documentclass[12pt]{article}

% =========================
% 0) Metadata (EDIT ME)
% =========================
\newcommand{\CourseCode}{MATH 521}
\newcommand{\CourseTitle}{Analysis I}
\newcommand{\CourseSection}{LEC 002}
\newcommand{\Semester}{Spring 2026}
\newcommand{\Instructor}{Thierry Laurens}
\newcommand{\MeetingInfo}{MWF 1:20--2:10 \;|\; VV B119} % optional
\newcommand{\HomeworkNumber}{7} % <-- change each week
\newcommand{\YourName}{Alejandro De La Torre}
\newcommand{\DueDate}{March 13, 2026} % <-- or set e.g. January 31, 2026

% =========================
% 1) Core math + proof tools
% =========================
\usepackage{amsmath, amssymb, amsthm}

% =========================
% 2) Lists and spacing
% =========================
\usepackage{enumitem}
\setlist[enumerate,1]{label=\textbf{(\alph*)}, leftmargin=2em, itemsep=0.25em}
\setlist[itemize]{leftmargin=1.5em, itemsep=0.25em}
\setlength{\parskip}{0.35em}
\setlength{\parindent}{0pt}

% =========================
% 3) Page geometry + header/footer
% =========================
\usepackage{geometry}
\geometry{margin=1in}

\usepackage{fancyhdr}
\pagestyle{fancy}
\fancyhf{}
\lhead{\CourseCode\ --- \CourseSection, \Semester \;(\MeetingInfo)}
\rhead{Homework \HomeworkNumber}
\cfoot{\thepage}
\renewcommand{\headrulewidth}{0.4pt}
% Fixes common fancyhdr warning about header height:
\setlength{\headheight}{15pt}
% Optional: reclaim a bit of vertical space after increasing headheight
\addtolength{\topmargin}{-2pt}

% =========================
% 4) Links (subtle)
% =========================
\usepackage[hidelinks]{hyperref}
\setcounter{tocdepth}{1}

% =========================
% 5) Quotes / figures
% =========================
\usepackage{csquotes}
\usepackage{graphicx}
\graphicspath{{images/}} % put figures in ./images/

% =========================
% 6) Simple scratch box (no tcolorbox)
% =========================
% A lightweight environment for brief planning / scratch work.
% This avoids any tcolorbox dependencies.
\newenvironment{scratch}{%
  \par\smallskip
  \noindent\textbf{Discussion \& Scratch Work}\begin{quote}\small
}{%
  \end{quote}\smallskip
}

% =========================
% 7) Lightweight theorem-like environments (optional)
% =========================
\newtheorem*{claim}{Claim}
\newtheorem*{lemma}{Lemma}
\newtheorem*{remark}{Remark}
\newtheorem{theorem}{Theorem}
\newtheorem{proposition}{Proposition}
\newtheorem{corollary}{Corollary}
\newtheorem{definition}{Definition}
\newtheorem{example}{Example}
\newtheorem{exercise}{Exercise}

% =========================
% 8) Problem header command (with TOC + PDF bookmarks)
% Usage: \problem[Optional short title]{<number>}
% =========================
\makeatletter
\newcommand{\problem}[2][]{%
  \def\prob@title{Problem #2\if\relax\detokenize{#1}\relax\else\ (\textit{#1})\fi}%
  \phantomsection%
  \pdfbookmark[1]{\prob@title}{prob#2}%
  \addcontentsline{toc}{section}{\prob@title}%
  \section*{\prob@title}%
  \vspace{-0.25em}\hrule\vspace{0.8em}%
}
\makeatother

% =========================
% 9) Title block
% =========================
\title{Homework \HomeworkNumber}
\author{\YourName \\
\CourseCode\ --- \CourseTitle \\
\CourseSection \\
Instructor: \Instructor}
\date{\DueDate}

\begin{document}
\maketitle

% ------------------ collaboration + sources ------------------
% \section*{Collaboration Statement}
% Write “I worked alone” or list collaborators / external resources.
% "I worked on this homework independently. No outside collaborators or resources were used."

\tableofcontents
\bigskip
\hrule
\bigskip

\newpage

% ============================================================
% Problems (duplicate blocks as needed)
% ============================================================

\problem[]{1}
For each of the following series, determine if it converges and prove your
answer.
\begin{enumerate}[label=\textbf{(\alph*)}]
\item \[
\sum_{n=2}^{\infty} \frac{1}{\log n}
\]
\item \[
\sum_{n=1}^{\infty} \frac{n-1}{n^2}
\]
\item \[
\sum_{n=1}^{\infty} \frac{n^2}{n!}
\]
\end{enumerate}

\begin{scratch}
% Outline the strategy, key definitions, lemmas, or a proof sketch.
\end{scratch}

\textbf{Proof.}\\
% Write your proof here.
\qed

\newpage

\problem[]{2}
For each of the following series, determine if it converges and prove your
answer.
\begin{enumerate}[label=\textbf{(\alph*)}]
\item \[
\sum_{n=1}^{\infty} \frac{n^4}{4^n}
\]
\item \[
\sum_{n=1}^{\infty} \frac{n!}{n^4+3}
\]
\item \[
\sum_{n=1}^{\infty} \frac{2^n}{n!}
\]
\end{enumerate}

\begin{scratch}
\end{scratch}

\textbf{Proof.}\\
\qed

\newpage

\problem[]{3}
For each of the following series, determine if it converges and prove your
answer.
\begin{enumerate}[label=\textbf{(\alph*)}]
\item \[
\sum_{n=1}^{\infty} \frac{1}{n^n}
\]
\item \[
\sum_{n=1}^{\infty} \frac{n!}{n^n}
\]
\item \[
\sum_{n=1}^{\infty} \frac{(-1)^n n!}{2^n}
\]
\end{enumerate}

\begin{scratch}
\end{scratch}

\textbf{Proof.}\\
\qed

\newpage

\problem[]{4}
Consider the series
\[
\sum_{n=1}^{\infty} 2^{(-1)^n-n}
= \frac{1}{4} + \frac{1}{2} + \frac{1}{16} + \frac{1}{8} + \frac{1}{64}
+ \frac{1}{32} + \cdots.
\]
\begin{enumerate}[label=\textbf{(\alph*)}]
\item Show that the ratio test is inconclusive.
\item Prove that the series converges using the root test.
\end{enumerate}

\begin{scratch}
\end{scratch}

\textbf{Proof.}\\
\qed

\newpage

\problem[]{5}
Below are two series $\sum_{n=1}^{\infty} a_n$ and the numbers $s\in\mathbb{R}$
to which they converge. For each of these series, what is the set of all
numbers $\tilde{s}\in\mathbb{R}$ such that there exists a rearrangement that
satisfies $\sum_{n=1}^{\infty} \tilde{a}_n=\tilde{s}$? Prove your answer.
\begin{enumerate}[label=\textbf{(\alph*)}]
\item \[
\sum_{n=1}^{\infty} \frac{(-1)^n}{n} = -\log 2
\]
\item \[
\sum_{n=1}^{\infty} \frac{(-1)^n}{n^2} = -\frac{\pi^2}{12}
\]
\end{enumerate}
(Hint: Your proof should be quite short, because we already did most of the work
in lecture.)

\begin{scratch}
\end{scratch}

\textbf{Proof.}\\
\qed

\newpage

\problem[]{6}
Prove or disprove the following statements:
\begin{enumerate}[label=\textbf{(\alph*)}]
\item If $\sum_{n=1}^{\infty} a_n$ converges, then $\sum_{n=1}^{\infty} a_{2n}$
converges.
\item If $\sum_{n=1}^{\infty} a_n$ converges absolutely, then
$\sum_{n=1}^{\infty} a_{2n}$ converges absolutely.
\end{enumerate}

\begin{scratch}
\end{scratch}

\textbf{Proof.}\\
\qed

% ============================================================
% Optional appendix: keep a personal toolbox of definitions/theorems
% ============================================================
\newpage
\appendix
\section*{Appendix: Toolbox}
\addcontentsline{toc}{section}{Appendix: Toolbox}

% Example (using amsthm environments instead of boxed tcolorboxes):
% \begin{definition}
% A metric space is ...
% \end{definition}
%
% \begin{theorem}[Heine--Borel]
% ...
% \end{theorem}

\end{document}
